\documentclass[a4paper,11pt]{article}
\usepackage[utf8]{inputenc}
\usepackage[margin=2.5cm]{geometry}
\usepackage{booktabs}
\usepackage{amsmath}
\usepackage{hyperref}

\title{HPPS Assignment 4: Locality Optimizations}
\author{Bastian Thomsen}
\date{December 2025}

\begin{document}

\maketitle

\section{Introduction}

This report documents my implementation of four programs for querying the OpenStreetMap place names dataset. The programs implement ID-based lookups (\texttt{id\_query\_naive}, \texttt{id\_query\_indexed}, \texttt{id\_query\_binsort}) and coordinate-based lookups (\texttt{coord\_query\_naive}).

I am confident in the overall quality of my solution. All programs compile without warnings (except for one in the provided \texttt{timing.h} which I cannot modify), produce correct results verified against each other, and have been tested with AddressSanitizer for memory safety.

To run the programs, use:
\begin{verbatim}
make all
./id_query_naive dataset.tsv
./id_query_indexed dataset.tsv
./id_query_binsort dataset.tsv
./coord_query_naive dataset.tsv
\end{verbatim}

\section{Analysis}

\subsection{Temporal and Spatial Locality}

\textbf{id\_query\_naive:} This implementation iterates through the full \texttt{struct record} array. Each \texttt{struct record} is large (containing many string pointers and numeric fields), so when scanning, we access much more data than needed (we only need \texttt{osm\_id}). \emph{Spatial locality is poor} because the relevant data (\texttt{osm\_id}) is spread far apart in memory. \emph{Temporal locality is minimal} since each record is visited once per query with no reuse.

\textbf{id\_query\_indexed:} This implementation creates a compact index of \texttt{struct index\_record} (16 bytes: \texttt{int64\_t} + pointer). When scanning, we access contiguous memory containing only the data we need. \emph{Spatial locality is significantly improved} because multiple \texttt{osm\_id} values fit in a single cache line (64 bytes can hold 4 index records). \emph{Temporal locality remains minimal} (single pass per query).

\textbf{id\_query\_binsort:} Uses the same compact index as \texttt{id\_query\_indexed}, but with binary search. \emph{Spatial locality is mixed}: the algorithm jumps around in the array, so consecutive accesses are not spatially local. However, the compact representation means each access loads useful data. \emph{Temporal locality is better for repeated queries}: the middle elements of the array are accessed frequently across queries, staying in cache.

\textbf{coord\_query\_naive:} Scans all records checking \texttt{lon} and \texttt{lat} fields. Similar to \texttt{id\_query\_naive}, \emph{spatial locality is poor} because we access a large struct for just two \texttt{double} values. \emph{Temporal locality is minimal}.

\subsection{Memory Corruption and Leaks}

I tested all programs using AddressSanitizer (compiled with \texttt{-fsanitize=address}) which detects memory errors including buffer overflows, use-after-free, and other corruption. No errors were reported.

For memory leaks, I verified that:
\begin{itemize}
    \item Each \texttt{malloc} has a corresponding \texttt{free}
    \item The index structures are freed via \texttt{free\_naive} or \texttt{free\_indexed}
    \item The record array is freed via \texttt{free\_records} (provided code)
    \item No memory is allocated without being freed
\end{itemize}

The code follows a simple allocation pattern: \texttt{mk\_*} allocates, \texttt{free\_*} deallocates. I do not store pointers to dynamically allocated memory elsewhere, so there are no dangling references.

\subsection{Correctness Confidence}

I am highly confident in the correctness of my implementations:

\begin{enumerate}
    \item All three ID query programs produce identical results when tested with 100 random valid IDs
    \item Results match the expected format from the assignment specification
    \item Edge cases tested: IDs that exist, IDs that do not exist (return ``not found'')
    \item The coordinate query correctly finds the closest place using Euclidean distance
\end{enumerate}

The implementations are straightforward: linear search, linear search on compact index, and binary search on sorted index. The algorithms are well-understood and the code closely follows standard patterns.

\subsection{Benchmarks}

I benchmarked the programs with different dataset sizes and query counts. Table~\ref{tab:benchmarks} shows results for 1000 queries.

\begin{table}[h]
\centering
\begin{tabular}{lrrrrr}
\toprule
Program & Records & Build (ms) & Query (ms) & Query/rec ($\mu$s) \\
\midrule
id\_query\_naive & 20,000 & 0 & 25.4 & 1.27 \\
id\_query\_indexed & 20,000 & 0 & 13.7 & 0.69 \\
id\_query\_binsort & 20,000 & 2 & 0.19 & 0.01 \\
\midrule
id\_query\_naive & 100,000 & 0 & 210.6 & 2.11 \\
id\_query\_indexed & 100,000 & 1 & 63.8 & 0.64 \\
id\_query\_binsort & 100,000 & 13 & 0.18 & 0.002 \\
\bottomrule
\end{tabular}
\caption{Benchmark results for 1000 ID queries}
\label{tab:benchmarks}
\end{table}

\textbf{Key observations:}

\begin{itemize}
    \item \textbf{Naive vs Indexed (linear search):} The indexed version is approximately 2x faster on 20k records and 3x faster on 100k records. This speedup comes from improved spatial locality: scanning 16-byte index records is faster than scanning large \texttt{struct record} entries.

    \item \textbf{Binary search:} The \texttt{id\_query\_binsort} is dramatically faster for queries ($O(\log n)$ vs $O(n)$). With 100k records, binary search is over 1000x faster per query. The trade-off is a build time of 13ms for sorting.

    \item \textbf{Amortization:} Binary search's build cost is amortized over queries. With 1000 queries, the 13ms build time adds only 13$\mu$s per query, which is negligible compared to the 210$\mu$s saved per query versus naive.

    \item \textbf{Scaling:} Query time for naive/indexed grows linearly with dataset size (5x records $\rightarrow$ ~5x time). Binary search query time remains nearly constant due to logarithmic complexity.
\end{itemize}

\textbf{Workload selection:} I chose 20k and 100k records to observe scaling behavior while keeping tests fast. I used 1000 queries to get stable timing averages and to demonstrate amortization of index build costs.

\end{document}
